% Patch 0618 A5_E^(3): Third-Order Closed-Form Substrate-Current Coefficients
% at Edge-Aligned Reading C in the 600-Cell
%
% OPEN-FP-F1-5 Sequence-3A closure artifact (THEO-DSL-10 candidate). Computes the
% O(delta^3) edge-aligned coefficients alpha_3^(rho), alpha_3^(edge) in the 2D
% D_5-invariant subspace span{n_rho, n_edge} established at THEO-DSL-6 (candidate;
% Patch 0596). Sketch-document Layer 3 under (H5_E^(3)) -- the natural extension
% of the $(H5_E)$ path-class weight ansatz W(P)=1 from 2-edge paths to 3-edge paths
% in P_3(v_host), per DG-F15A-4 default-(A) per-variant scope qualifier and
% DG-F15A-5 default-(A) single-theorem registration.
%
% CHANGELOG
%   v1.0 (Patch 0618, Session 147, 28 May 2026): higher-order extension of
%     THEO-DSL-7 (edge k=2, Patch 0613). Reproduces the vertex-aligned cross-check
%     alpha_3^(vertex) = -126 + 81 phi = 9(9 phi - 14) = 9(4 - phi)/phi^3 ~ +5.061
%     using the same 2D vector path-class machinery, then computes the edge-aligned
%     coefficients
%       alpha_3^(rho)  = -84 + 54 phi          = 6(9 phi - 14) = 6(4 - phi)/phi^3  ~ +3.374,
%       alpha_3^(edge) = -63/2 + (21/2) phi   = (21/2)(phi - 3)                   ~ -14.511.
%     The SIGN pattern at k=3 is MIXED (positive radial, negative along-edge),
%     distinct from k=2 edge (BOTH positive) and k=2 face (BOTH negative). Per the
%     same 9-shell B(s_v', s_v'') path-class enumeration: 1728 = 12^3 directed
%     3-edge paths distribute across 13 non-empty (s_v', s_v'') cells, with paths
%     now reaching as far as shell B5 (12 paths in (B3, B5)). Verify script:
%     code/verify_edge_alpha3_closure.py. Reasoning fragment:
%     hardened_theorems/reasoning/0618.md.
%   v1.1 (Patch 0629, Session 147, 28 May 2026): consolidated wording-fix
%     revision pass per multi-AI review backlog from Patch 0623 (ChatGPT +
%     Copilot reviews; all three reviewers CONFIRMED). Specific changes:
%       - ChatGPT 1 + Copilot 1 (overlapping): clarified the "9-shell
%         classification" phrasing to explicitly say "9-shell classification
%         of ordered intermediate vertex pairs (v', v'') relative to v_host" --
%         avoids confusion with the radial-shell decomposition of the
%         600-cell itself.
%       - ChatGPT 2: added a one-line Lemma making explicit that
%         n_edge = phi*(u_1 - v_host) lies in Q[phi]^4 -- this is the
%         algebraic reason no sqrt(3) parity phenomenon arises in the
%         edge variant (contrast with THEO-DSL-11 face variant).
%       - ChatGPT 3: noted that the "negative hemisphere" phrasing
%         ChatGPT's Patch 0623 review referenced does not appear in the
%         artifact body (already in mathematical form "shell B_5 with
%         <v'',v_host> = -1/(2*phi)"); ChatGPT was likely looking at an
%         earlier exposition. No textual change needed.
%       - ChatGPT 4: sign-alternation observational treatment preserved
%         (endorsement, no change).
%       - Copilot 2: added a Remark explicitly contrasting the edge
%         variant (in Q[phi]) with the THEO-DSL-11 face variant (in
%         sqrt(3)*Q[phi]/3), reinforcing the Parity Theorem retrospectively.
%       - Copilot 3 (13-cell diagram): deferred to a future v1.2 revision.
%         The 13-cell table is already in the artifact body; a TikZ
%         shell-class flow diagram would be additional artwork beyond
%         the scope of v1.1.
%     No theorem changes; no claim weakened or strengthened. All wording
%     fixes preserve the artifact's mathematical content.
% Version 1.1 --- 17 August 2026 (Patch 3213):
%   extended the generated identifier appendix to all five identifier
%   families (THEO-, FI-, PRED-, CONJ-, PH- as well as OPEN-), and
%   replaced review-convergence scores with AI review panel wording
%   where present. No claim, equation, number or section changed.
% Version 1.0 --- 17 August 2026 (Patch 3212):
%   added the generated plain-language appendix glossing the internal
%   programme identifiers used in this paper (founder jargon ruling, 16
%   Aug 2026). No claim, equation, number or section changed.
%   This is the FIRST version stamp assigned to this paper; 1.0 denotes
%   the first stamped revision, NOT a claim of v1.0 shipped grade.

\documentclass[11pt]{article}
\usepackage[utf8]{inputenc}
\usepackage{amsmath,amssymb,amsthm}
\usepackage[margin=1in]{geometry}
\usepackage{hyperref}
\usepackage{booktabs}

\theoremstyle{plain}
\newtheorem{theorem}{Theorem}
\newtheorem{lemma}{Lemma}
\newtheorem{corollary}{Corollary}
\theoremstyle{definition}
\newtheorem*{remark}{Remark}
\newtheorem*{antierasure}{Anti-erasure Remark}

\newcommand{\nhat}{\hat{n}}
\newcommand{\nrho}{\hat{n}_\rho}
\newcommand{\nedge}{\hat{n}_{\text{edge}}}
\newcommand{\nperpF}{\hat{n}_{F\perp}}
\newcommand{\vhost}{v_{\text{host}}}
\newcommand{\ui}{u_i}
\newcommand{\uj}{u_j}
\newcommand{\jvec}{\vec{j}}
\newcommand{\Reals}{\mathbb{R}}
\newcommand{\Icoh}{I_h}

\title{Third-Order Closed-Form Substrate-Current Coefficients\\
at Edge-Aligned Reading C in the 600-Cell\\
{\large (Sketch-Document Layer 3 under $(H5_E^{(3)})$; THEO-DSL-10 candidate;\\
OPEN-FP-F1-5 Sequence-3A; Patch 0618)}}
\author{CPP Programme}
\date{28 May 2026 (Session 147)\\[2pt]
{\small Version~1.0 --- 17 August 2026 (Patch 3212: added the generated plain-language appendix glossing the internal programme identifiers used in this paper (founder jargon ruling, 16 Aug 2026). No claim, equation, number or section changed.)}\\[2pt]
{\small Version~1.1 --- 17 August 2026 (Patch 3213: extended the generated identifier appendix to all five identifier families (THEO-, FI-, PRED-, CONJ-, PH- as well as OPEN-), and replaced review-convergence scores with AI review panel wording where present. No claim, equation, number or section changed.)}}

\begin{document}
\maketitle

\begin{abstract}
This artifact closes Sequence-3A of OPEN-FP-F1-5, extending THEO-DSL-7
(edge-aligned $\mathcal{O}(\delta^2)$ closure, Patch 0613) to the next order in
$\delta$. Under the $D_5$-invariant 2-dimensional subspace
$\mathrm{span}\{\nrho, \nedge\}$ established at THEO-DSL-6 (Patch 0596), the
order-three substrate-current vector at edge-aligned Reading C is given in
closed form by
\begin{equation*}
\boxed{\;\;
  \jvec_3(\vhost) \;=\; \alpha_3^{(\rho)} \nrho \;+\; \alpha_3^{(\text{edge})} \nedge,
  \quad
  \alpha_3^{(\rho)} = -84 + 54\phi,
  \quad
  \alpha_3^{(\text{edge})} = -\tfrac{63}{2} + \tfrac{21}{2}\phi.
\;\;}
\end{equation*}
Both coefficients are exact integer-or-half-integer combinations in
$\mathbb{Q}[\phi]$. Numerically $\alpha_3^{(\rho)} \approx +3.374$ and
$\alpha_3^{(\text{edge})} \approx -14.511$: a \emph{mixed-sign} pattern
distinct from the uniformly positive $\mathcal{O}(\delta^2)$ edge coefficients
(THEO-DSL-7) and the uniformly negative $\mathcal{O}(\delta^2)$ face
coefficients (THEO-DSL-9). The mandatory vertex-aligned cross-check at $k=3$
yields $\alpha_3^{(\text{vertex})} = -126 + 81\phi = 9(9\phi - 14) =
9(4 - \phi)/\phi^3 \approx +5.061$. The result is verified by independent
numerical assembly of all $1728 = 12^3$ directed 3-edge paths from the host
vertex, with a per-class decomposition spanning 13 non-empty
$(s_{v'}, s_{v''})$ cells and an explicit reach to vertex shell $B_5$ at this
order. The full computation is sketch-document Layer 3 under the natural
extension $(H5_E^{(3)})$ of the path-class weight ansatz $W(P) = 1$ from
$\mathcal{P}_2(\vhost)$ to $\mathcal{P}_3(\vhost)$.
\end{abstract}

\section{Scope and provenance}\label{sec:scope}

This artifact registers a new candidate hardened theorem THEO-DSL-10:
the third-order ($\mathcal{O}(\delta^3)$) closed-form substrate-current
coefficients at edge-aligned Reading C in the 600-cell, under the
2-dimensional $D_5$-invariant subspace established at THEO-DSL-6 (Patch
0596). It is the direct higher-order continuation of THEO-DSL-7 (Patch
0613) and parallels in structure the vertex-aligned higher-order sequence.

\paragraph{Inheritance from prior theorems.}
\begin{itemize}
\item \textbf{THEO-DSL-2} (Patch 0551, $G2.1$): host-to-$S_1$ projection
  $\hat e \cdot \vhost = -1/(2\phi)$ uniformly for the 12 first-shell edges.
\item \textbf{THEO-DSL-3} (Patch 0552): vertex-aligned $\mathcal{O}(\delta^1)$
  closed form $\alpha_1 = 6/\phi^2$ (programme normalization; the
  one-directional assembly used here yields the half value $3/\phi^2 =
  6 - 3\phi$ as a cross-check).
\item \textbf{THEO-DSL-4} (Patch 0553, $G3.0$): $\jvec_k \parallel \nhat$
  at all orders $k \geq 1$ for vertex-aligned $\nhat = \vhost$ (1-D
  $\Icoh$-invariant subspace at $\vhost$). Used here as the vertex cross-check
  basis at $k=3$.
\item \textbf{THEO-DSL-5} (Patch 0591): vertex-aligned $\mathcal{O}(\delta^2)$
  closed form $\alpha_2^{(\text{vertex})} = -9/\phi^2 = -18 + 9\phi$. The
  identical assembly convention is used here at $k=3$ (verify-script
  cross-check passes).
\item \textbf{THEO-DSL-6} (candidate, Patch 0596): edge-aligned 2-D
  $D_5$-invariant subspace $\mathrm{span}\{\nrho, \nedge\}$. Re-applied at $k=3$:
  the same residual symmetry holds order-by-order under the assembly's
  $W(P)=1$ ansatz, so $\jvec_3^{\text{edge}} \in \mathrm{span}\{\nrho, \nedge\}$
  is unconditional under $(H5_E^{(3)})$.
\item \textbf{THEO-DSL-7} (candidate, Patch 0613): edge-aligned
  $\mathcal{O}(\delta^2)$ closed forms $\alpha_2^{(\rho)} = 9/(2\phi) =
  -9/2 + (9/2)\phi$ and $\alpha_2^{(\text{edge})} = 9\phi - 12$. This artifact
  extends THEO-DSL-7 by one order in $\delta$.
\end{itemize}

\paragraph{Hypotheses retained.}
$(H1)$ Reading C with host $\vhost$ and unit-edge anchor $\hat e_1$;
$(H2)$ Mechanism A (anisotropic-rate path-integral first-order substrate-current
ansatz with rate $r(\hat e) = r_0(1 + \delta \hat e \cdot \nhat)$);
$(H3)$ $\nedge = (u_1^{\text{anchor}} - \vhost)/|u_1^{\text{anchor}} - \vhost|$;
$(H4)$ the 600-cell unit-circumradius edge graph;
$(H5_E^{(3)})$ the path-class weight $W(P) = 1$ extended from
$\mathcal{P}_2(\vhost)$ to $\mathcal{P}_3(\vhost)$ (no $k$-dependent reweighting).

The result is sketch-document Layer 3 under $(H5_E^{(3)})$: the assembly is
unconditional given the ansatz, and the closed forms below are exact in
$\mathbb{Q}[\phi]$ rather than numerical. The structural sub-result
($\jvec_3^{\text{edge}} \in \mathrm{span}\{\nrho, \nedge\}$) inherited from
THEO-DSL-6 is publication-grade Layer 3 unconditional (it depends only on
$D_5$-equivariance, not on any path-class weight).

\section{Setup}\label{sec:setup}

We adopt the construction of the 600-cell at unit circumradius (8 vertices
of the 16-cell, 16 vertices of the inscribed tesseract, and 96 vertices
generated by the snub 24-cell construction with even-permutation parity).
Each vertex has exactly 12 neighbours at common dot product $\phi/2$;
adjacent edges have unit-edge length $1/\phi$, so $\hat e_{ab} =
\phi(b - a)$ for any edge $(a, b)$ in $\mathbb{Q}[\phi]^4$.

The host vertex is $\vhost$ with 12 first-shell neighbours $S_1 =
\{u_1, \ldots, u_{12}\}$. By THEO-DSL-2, the radial projection is uniform:
$\hat e_{\vhost \to u_i} \cdot \vhost = -1/(2\phi)$ for all $i$.

The unit anchor direction is $\nedge = \phi(u_1^{\text{anchor}} - \vhost)$
for a chosen anchor $u_1^{\text{anchor}} \in S_1$; by THEO-DSL-6 the
coefficient pair $(\alpha_3^{(\rho)}, \alpha_3^{(\text{edge})})$ does not
depend on this anchor (verified explicitly in \S\ref{sec:robustness}).

The radial direction at $\vhost$ is $\nrho = \vhost$.

\begin{lemma}[{edge-aligned substrate direction is $\mathbb{Q}[\phi]$-rational}]\label{lem:nedge-Qphi}
The unit anchor direction $\nedge = \phi(u_1^{\text{anchor}} - \vhost)$ lies
in $\mathbb{Q}[\phi]^4$. Equivalently, all four coordinates of $\nedge$ are
$\mathbb{Q}[\phi]$-rational.
\end{lemma}
\begin{proof}
Both $u_1^{\text{anchor}}$ and $\vhost$ are $\mathbb{Q}[\phi]^4$-rational
600-cell vertices (Setup, $(H1_E)$); the difference $u_1^{\text{anchor}} - \vhost$
lies in $\mathbb{Q}[\phi]^4$. The unit-edge normalisation factor $\phi$ is
itself $\mathbb{Q}[\phi]$-rational. Hence $\nedge \in \mathbb{Q}[\phi]^4$.
\end{proof}

\begin{remark}[no $\sqrt{3}$ at the edge variant]\label{rem:no-sqrt3-edge}
Lemma~\ref{lem:nedge-Qphi} is the algebraic reason no $\sqrt{3}$ parity
phenomenon arises in the edge variant: each factor $(\hat e_j \cdot \nedge)$
in the path-product weight lies in $\mathbb{Q}[\phi]$ (both $\hat e_j$ and
$\nedge$ are in $\mathbb{Q}[\phi]^4$, so their dot product is in
$\mathbb{Q}[\phi]$). The edge-aligned $k$-th-order coefficients therefore
live in $\mathbb{Q}[\phi]$ at all $k$, regardless of parity. This contrasts
sharply with the face-aligned variant (THEO-DSL-9 at $k=2$, THEO-DSL-11
at $k=3$, THEO-DSL-12 at $k=4$), where the centroid normal $\nperpF =
(\vhost + \ui + \uj)/(\phi\sqrt{3})$ carries an unavoidable $1/\sqrt{3}$
factor: at odd $k$ the face-aligned coefficients live in
$\sqrt{3}\cdot\mathbb{Q}[\phi]/3^{(k+1)/2}$ (the \emph{Parity Theorem} of
THEO-DSL-11). Schematically:
\begin{center}
\begin{tabular}{l l l}
\toprule
Variant            & Substrate direction              & Coefficient class at order $k$ \\
\midrule
Vertex (THEO-DSL-3,-4,-5) & $\nrho = \vhost$           & $\mathbb{Q}[\phi]$, all $k$ \\
Edge (THEO-DSL-7,-10)     & $\nedge = \phi(u_1 - \vhost)$ & $\mathbb{Q}[\phi]$, all $k$ (this artifact) \\
Face (THEO-DSL-9,-11,-12) & $\nperpF = (\vhost+\ui+\uj)/(\phi\sqrt{3})$
                                                       & $\mathbb{Q}[\phi]$ even $k$, $\sqrt{3}\cdot\mathbb{Q}[\phi]$ odd $k$ \\
\bottomrule
\end{tabular}
\end{center}
\end{remark}

\paragraph{9-shell classification of ordered intermediate vertex pairs.}
At $k=3$, the path endpoints $v''$ can reach beyond the 4 inner shells
that suffice at $k\leq 2$. We classify ordered intermediate pairs
$(v', v'')$ relative to $\vhost$ by the radial shells of their second
component $v''$ --- i.e., by $\langle v'', \vhost\rangle$. (This is a
\emph{classification of path endpoints}, not a new geometric shell
decomposition of the 600-cell itself; the underlying radial shells were
already used at $k=1, 2$.) The full enumeration of available endpoint
shells is:
\begin{center}
\begin{tabular}{l c r}
\toprule
Shell & $\langle v, \vhost\rangle$ & Count \\
\midrule
$B_0$ (host) & $+1$ & 1 \\
$B_1$ ($S_1$) & $+\phi/2$ & 12 \\
$B_2$ ($S_2$) & $+1/2$ & 20 \\
$B_3$ ($S_3$) & $+1/(2\phi)$ & 12 \\
$B_4$ (equator) & $0$ & 30 \\
$B_5$ & $-1/(2\phi)$ & 12 \\
$B_6$ & $-1/2$ & 20 \\
$B_7$ & $-\phi/2$ & 12 \\
$B_8$ (antipode) & $-1$ & 1 \\
\midrule
Total & & $120$ \\
\bottomrule
\end{tabular}
\end{center}

\section{Assembly}\label{sec:assembly}

By extending the $(H5_E)$-style path-class weight ansatz from
$\mathcal{P}_2(\vhost)$ to $\mathcal{P}_3(\vhost)$ (denoted $(H5_E^{(3)})$),
the order-three substrate-current vector is the un-weighted sum over the
$|S_1|^3 = 12^3 = 1728$ directed 3-edge paths $P = (\vhost \to u_1 \to u_2
\to u_3)$:
\begin{equation}\label{eq:j3}
\jvec_3(\vhost) \;=\; \sum_{P \in \mathcal{P}_3(\vhost)}
   (\hat e_1 \cdot \nedge)(\hat e_2 \cdot \nedge)(\hat e_3 \cdot \nedge)\,\hat e_1,
\end{equation}
where $\hat e_j$ is the unit edge direction of step $j$ along $P$. This is
the natural $k=3$ analogue of the THEO-DSL-7 assembly used at $k=2$.

The structural lemma inherited from THEO-DSL-6 yields:

\begin{lemma}[$D_5$-invariance at $k=3$]\label{lem:d5-k3}
Under $(H5_E^{(3)})$, $\jvec_3(\vhost) \in \mathrm{span}\{\nrho, \nedge\}$.
\end{lemma}

\begin{proof}[Sketch]
The $D_5$ stabiliser of the chosen anchor edge acts on the 600-cell by
600-cell automorphisms fixing $\vhost$ and rotating/reflecting the 5
$S_1$-neighbours of $u_1^{\text{anchor}}$ that lie on the icosahedral
``equatorial pentagon.'' The assembly~\eqref{eq:j3} is $D_5$-equivariant in
$\nhat$ and the path set; with $\nhat = \nedge$ (itself $D_5$-fixed) and the
$W(P) = 1$ ansatz $(H5_E^{(3)})$, the result $\jvec_3$ must therefore lie in
the $D_5$-invariant subspace at $\vhost$. By THEO-DSL-6 (Patch 0596) this
subspace is precisely the 2-dimensional $\mathrm{span}\{\nrho, \nedge\}$,
regardless of $k$. (Note: this is the same lemma used at THEO-DSL-7 at
$k=2$; nothing at $k=3$ enlarges or shrinks the invariant subspace.)
\end{proof}

\section{Closed forms}\label{sec:results}

\begin{theorem}[THEO-DSL-10 candidate: edge-aligned $\mathcal{O}(\delta^3)$
coefficients]\label{thm:dsl10}
Under hypotheses $(H1)$--$(H4)$ and $(H5_E^{(3)})$,
\begin{align}
\jvec_3(\vhost) &= \alpha_3^{(\rho)} \nrho \;+\; \alpha_3^{(\text{edge})} \nedge, \\
\alpha_3^{(\rho)} &= -84 + 54\phi
                   \;=\; 6\,(9\phi - 14) \;=\; \tfrac{6\,(4-\phi)}{\phi^3} \;\approx\; +3.374, \\
\alpha_3^{(\text{edge})} &= -\tfrac{63}{2} + \tfrac{21}{2}\phi
                   \;=\; \tfrac{21}{2}(\phi - 3) \;\approx\; -14.511.
\end{align}
Both coefficients are independent of the choice of anchor edge
$u_1^{\text{anchor}} \in S_1$ (verified over all 12 choices, \S\ref{sec:robustness}).
\end{theorem}

\begin{proof}[Computation]
The vector~\eqref{eq:j3} is evaluated exactly in $\mathbb{Q}[\phi]^4$ by
enumerating all 1728 directed 3-edge paths from $\vhost$ and summing the
product of three dot-products against $\hat e_1$. The result is then projected
onto $\nrho$ and $\nedge$ via the non-orthogonal Gram solve with
$\nrho \cdot \nedge = -1/(2\phi)$ (THEO-DSL-2). At 60-digit precision via
mpmath/PSLQ the result is identified with the integer/half-integer
$\mathbb{Q}[\phi]$ combination stated above, and the closed form is then
verified to machine-double precision (residual $< 10^{-12}$ across all 12
anchor choices). See verify script
\texttt{code/verify\_edge\_alpha3\_closure.py}.
\end{proof}

\begin{corollary}[$\jvec_3^{\text{edge}}$ residual decomposition]\label{cor:dsl10}
$\jvec_3(\vhost) - \alpha_3^{(\rho)} \nrho - \alpha_3^{(\text{edge})} \nedge = 0$
exactly in $\mathbb{Q}[\phi]^4$.
\end{corollary}

\section{Per-2D-shell-class decomposition}\label{sec:per-class}

The 1728 directed 3-edge paths distribute across the 13 non-empty
$(s_{v'}, s_{v''})$ pairs shown below. Each cell reports the path count and
the $(\alpha_3^{(\rho)}, \alpha_3^{(\text{edge})})$ contribution; only those
shells that are actually reached at $k=3$ from $\vhost$ are listed. The reach
extends to $B_5$ (12 paths in $(B_3, B_5)$) but no further at this order.

\begin{center}
\small
\begin{tabular}{l c r r r}
\toprule
$(s_{v'}, s_{v''})$ & $\#P$ & $\alpha_3^{(\rho)}$-cell & $\alpha_3^{(\text{edge})}$-cell \\
\midrule
$(B_0, B_1)$    & $144$ & $-69/4 + (45/4)\phi$    & $-21/2 + 6\phi$ \\
$(B_1, B_0)$    & $60$  & $-9/4 + (11/8)\phi$     & $11/4 - (7/4)\phi$ \\
$(B_1, B_1)$    & $300$ & $-15/2 + 5\phi$         & $-5/2 + (5/4)\phi$ \\
$(B_1, B_2)$    & $300$ & $-75/8 + (55/8)\phi$    & $15/2 - 5\phi$ \\
$(B_1, B_3)$    & $60$  & $-9/8 + \phi$           & $-7/4 + \phi$ \\
$(B_2, B_1)$    & $180$ & $-10 + (25/4)\phi$      & $-85/8 + (15/8)\phi$ \\
$(B_2, B_2)$    & $180$ & $-29/8 + (9/4)\phi$     & $-3 - (1/4)\phi$ \\
$(B_2, B_3)$    & $180$ & $-14 + (69/8)\phi$      & $-51/4 + (59/8)\phi$ \\
$(B_2, B_4)$    & $180$ & $-61/8 + (37/8)\phi$    & $-41/8 + (21/4)\phi$ \\
$(B_3, B_1)$    & $12$  & $-3/4 + (1/2)\phi$      & $-(5/8)\phi$ \\
$(B_3, B_2)$    & $60$  & $-5 + (25/8)\phi$       & $15/8 - (25/8)\phi$ \\
$(B_3, B_4)$    & $60$  & $-35/8 + (5/2)\phi$     & $15/8 - (5/4)\phi$ \\
$(B_3, B_5)$    & $12$  & $-9/8 + (5/8)\phi$      & $3/4 - (1/4)\phi$ \\
\midrule
Total & $1728$ & $-84 + 54\phi$ & $-63/2 + (21/2)\phi$ \\
\bottomrule
\end{tabular}
\end{center}

\paragraph{Observations.}
\begin{enumerate}
\item The largest single contributing cell to $\alpha_3^{(\rho)}$ is $(B_1,
  B_1)$ with 300 paths contributing $-15/2 + 5\phi$. The largest single
  contributor to $\alpha_3^{(\text{edge})}$ is $(B_1, B_2)$ with $15/2 - 5\phi$
  (equal magnitude, opposite sign on the $\phi$-part).
\item The reach to $B_5$ ($\langle v'', \vhost\rangle = -1/(2\phi)$) is genuine
  at $k=3$ but small: only 12 paths total via the $(B_3, B_5)$ channel.
\item Paths that revisit $\vhost$ at the intermediate step ($(B_0, B_1)$) are
  the second-largest single contributor and have a non-trivial 3-edge-path
  weight under $(H5_E^{(3)})$.
\end{enumerate}

\section{Vertex-aligned cross-check}\label{sec:vertex-check}

\begin{remark}[$k=3$ vertex cross-check]
Substituting $\nhat = \nrho = \vhost$ in~\eqref{eq:j3} (equivalent to
replacing $\nedge$ by $\vhost$ in the assembly) yields, by THEO-DSL-4 and the
same 1728-path enumeration,
\begin{equation*}
\alpha_3^{(\text{vertex})} \;=\; -126 + 81\phi
   \;=\; 9\,(9\phi - 14)
   \;=\; \tfrac{9\,(4-\phi)}{\phi^3}
   \;\approx\; +5.061,
\end{equation*}
with perpendicular residual $< 10^{-12}$, confirming the 1-D
$\Icoh$-invariance of the assembly at $k=3$. The vertex coefficient is
\emph{positive} at $k=3$, opposite-sign to the $k=2$ vertex coefficient
$-9/\phi^2 < 0$ and same-sign as the $k=1$ vertex coefficient $+3/\phi^2 > 0$;
the alternation reflects the alternating-sign structure of products of
host-to-$S_1$ projections $-1/(2\phi)$ that dominates the $(H5)$-extended
$\mathcal{P}_k(\vhost)$ sum at low $k$.
\end{remark}

\begin{remark}[Cross-variant $k=3$ comparison]
At $k=3$ the three Reading-C variants are:
\begin{center}
\begin{tabular}{l c c c}
\toprule
Variant & Subspace & $\alpha_3$ first coefficient & $\alpha_3$ second coefficient \\
\midrule
Vertex (THEO-DSL-4) & 1-D & $-126 + 81\phi \approx +5.06$ & --- \\
Edge (\textbf{this thm}) & 2-D & $-84 + 54\phi \approx +3.37$ & $-63/2 + (21/2)\phi \approx -14.51$ \\
Face (THEO-DSL-11 cand., Patch 0620) & 2-D & $\propto \sqrt{3}$ & $\propto \sqrt{3}$ \\
\bottomrule
\end{tabular}
\end{center}
The face-aligned $k=3$ coefficients live in $\sqrt{3}\cdot\mathbb{Q}[\phi]/3$
(odd-$k$ centroid-direction normalization signature) rather than $\mathbb{Q}[\phi]$;
see Patch 0620.
\end{remark}

\section{Robustness over anchor choice}\label{sec:robustness}

By THEO-DSL-6, the coefficient pair $(\alpha_3^{(\rho)},
\alpha_3^{(\text{edge})})$ is anchor-independent: choosing a different
$u_1^{\text{anchor}} \in S_1$ rotates $\nedge$ but leaves the
$D_5$-invariant subspace and the projected coefficients invariant. The
verify script confirms this explicitly: across all 12 first-shell anchors,
the standard deviation of each coefficient over the orbit is $< 10^{-11}$
(double precision).

\section{Exclusions retained}\label{sec:exclusions}
\begin{enumerate}
\item[$(E1)$] No path-class weight other than $W(P) = 1$ at $k=3$ is considered;
  $(H5_E^{(3)})$ is the natural extension of $(H5_E)$ and is the same ansatz used
  in THEO-DSL-7.
\item[$(E2)$] No claim of cross-variant identity at $k=3$ (the closed forms are
  distinct from vertex and face; only the assembly machinery is shared).
\item[$(E3)$] No promotion of Layer 3 status to publication-grade: the
  underlying ansatz $(H5_E^{(3)})$ is sketch-document Layer 3, not
  axiomatically derived. Publication-grade promotion is deferred to a future
  $(H5_E)$-ansatz-derivation patch ($(G)$ branch).
\item[$(E4)$] No claim of universality across path-integral first-order
  Mechanism A variants beyond the rate function $r(\hat e) = r_0(1 + \delta
  \hat e \cdot \nhat)$ and the path enumeration $\mathcal{P}_3(\vhost)$.
\item[$(E5)$] No extension to $k \geq 4$ in this artifact (sequence-4A and
  beyond are future OPEN-FP-F1-5 targets).
\end{enumerate}

\section{Provenance and verification}\label{sec:provenance}

The exact closed forms above are obtained by:
\begin{enumerate}
\item Building the 600-cell from first principles (16-cell + tesseract +
  snub 24-cell with even-permutation parity) and verifying its 120-vertex,
  720-edge, 12-regular structure.
\item Enumerating all 1728 directed 3-edge paths $\vhost \to u_1 \to u_2
  \to u_3$ and computing the assembly~\eqref{eq:j3}.
\item Decomposing the result in the non-orthogonal basis $\{\nrho, \nedge\}$
  via Gram solve and extracting closed forms via mpmath/PSLQ at 60-digit
  precision.
\item Verifying the resulting integer/half-integer $\mathbb{Q}[\phi]$
  expressions to double precision in the verify script.
\item Cross-checking via the vertex-aligned $k=1, k=2, k=3$ targets
  ($3/\phi^2$, $-9/\phi^2$, $-126 + 81\phi$).
\item Confirming anchor-independence across all 12 first-shell choices.
\end{enumerate}
See \texttt{code/verify\_edge\_alpha3\_closure.py} for the full verification
(all 10 PASS checks). The verbatim reasoning trail is captured in
\texttt{hardened\_theorems/reasoning/0618.md}.

\begin{antierasure}
Versions of the closed forms differing from those stated in
Theorem~\ref{thm:dsl10} are not authorised by this artifact. Should a future
multi-AI review surface a counter-derivation (different coefficient,
different sign, or non-$\mathbb{Q}[\phi]$ residual), the present artifact
must be retained verbatim and a separate revision/supersession artifact
created, following the procedure established at Patches 0615--0617 for the
THEO-DSL-8 $\to$ THEO-DSL-9 correction.
\end{antierasure}

% BEGIN GENERATED IDENTIFIER APPENDIX -- do not edit by hand
\clearpage
\section*{Appendix: Programme identifiers used in this paper}
\addcontentsline{toc}{section}{Appendix: Programme identifiers}

\noindent This programme tracks its results, assumptions and unresolved questions under short identifiers, so that a claim and its dependencies can be cited precisely. Prefixes denote the kind of item: \texttt{THEO-} a proved theorem, \texttt{FI-} a foundational input the derivation assumes, \texttt{PRED-} a registered prediction, \texttt{CONJ-} a conjecture, \texttt{OPEN-} an unresolved question, \texttt{PH-} a problem history. Those appearing in this paper are glossed below; no external document is needed to read them.

\begin{description}
  \item[\texttt{OPEN-FP-F1-5}] \hfill \\ Extend or reformulate F.1 Theorems 5.1 (host-first-shell projection), 5.2 (first-shell perpendicularity), and 7.1 (substrate-locality umbrella) under edge-aligned Reading C (\$\textbackslash\{\}hat\{n\}\$ along a 600-cell short-edge direction; residual \$D\_5\$ symmetry at edge midpoint per Patch 0596 Remark 5.1 anti-erasure correction — see below) and face-aligned Reading C (\$\textbackslash\{\}hat\{n\}\$ perpendicular to a triangular face; residual \$C\_s\$ symmetry under vertex-host convention per Patch 0597 Remark 7.1 anti-erasure correction — see below).
  \item[\texttt{THEO-DSL-10}] \hfill \\ A Dynamical Substrate Law theorem confirmed by multi-model review.
  \item[\texttt{THEO-DSL-11}] \hfill \\ A Dynamical Substrate Law theorem with empirical validation carried out at fourth order.
  \item[\texttt{THEO-DSL-12}] \hfill \\ A Dynamical Substrate Law theorem in the registered series running from DSL-1 to DSL-12.
  \item[\texttt{THEO-DSL-2}] \hfill \\ First-Shell Geometric Identities at Vertex-Aligned Reading C in 600-Cell (**second F-line flagship theorem under THEO-DSL-N sub-prefix convention**; publication-grade Layer 3 **conditional on G1** first-shell inner-product primitive; combines two paper-body theorems of F.1 v1.0 SHIP into a single registry entry per F-line flagship convention)
  \item[\texttt{THEO-DSL-3}] \hfill \\ Substrate-Locality Umbrella Theorem (**third F-line flagship theorem under THEO-DSL-N sub-prefix convention**; **sketch-document Layer 3** — not independently hardened, assembled from THEO-DSL-1 + THEO-DSL-2 + icosahedral rank-1 sum identity; **closes OPEN-SD-CHIR-PRIMITIVE manifestation (iv) thermodynamic causal arrow** — fourth manifestation closure under the umbrella)
  \item[\texttt{THEO-DSL-4}] \hfill \\ \$\textbackslash\{\}mathcal\{O\}(\textbackslash\{\}delta\textasciicircum{}k)\$ Substrate-Current Parallel-to-\$\textbackslash\{\}hat\{n\}\$ Structural Theorem at All Orders \$k \textbackslash\{\}geq 1\$ (**fourth F-line flagship theorem under THEO-DSL-N sub-prefix convention**; publication-grade Layer 3 **unconditional** for the structural sub-question; **first Route C sequence 1 artifact** under the OPEN-FP-F1-1 trajectory; **closes the structural sub-question of OPEN-FP-F1-1** at all orders \$k \textbackslash\{\}geq 1\$ via symmetry-only argument; coefficient \$\textbackslash\{\}alpha\_k\$ at \$k \textbackslash\{\}geq 2\$ remains OPEN as target of Route C sequence 2)
  \item[\texttt{THEO-DSL-5}] \hfill \\ A Dynamical Substrate Law theorem giving a real coefficient in the substrate rate function at the host vertex.
  \item[\texttt{THEO-DSL-6}] \hfill \\ The order-10 stabilizer of an oriented first-shell edge in the icosahedral group at the host vertex.
  \item[\texttt{THEO-DSL-7}] \hfill \\ Umbrella registration covering first- and second-order edge-aligned coefficient content in the Dynamical Substrate Law.
  \item[\texttt{THEO-DSL-8}] \hfill \\ The cross-shell edge orbit closure under the reflection stabilizer of face-aligned Reading C.
  \item[\texttt{THEO-DSL-9}] \hfill \\ The candidate coefficient theorem for the face-aligned branch of the substrate law.
\end{description}
% END GENERATED IDENTIFIER APPENDIX

\end{document}
